\chapter{Bulut ve Sanallastirma Guvenligi}

\section*{Giris}
Bulut ve sanallastirma guvenligi, paylasimli sorumluluk modeliyle ele alinmalidir. Kimlik, ag, veri ve is yukleri icin bulut-yerel kontrollerin bir butun olarak calismasi gerekir.

\section{Buyuk Olcekli Sanallastirma Mimarilerinde Guvenlik}
\subsection{Hipervizor Guvenligi: VMware, Proxmox ve KVM Ortamlarinin Sikilastirilmasi}
Hipervizor katmaninda yonetim arayuzu erisim kontrolu, ag izolasyonu ve host hardening temel savunmadir.

\section{Bulut Bilisim Servis Modelleri (IaaS, PaaS, SaaS)}
IaaS/PaaS/SaaS modellerinde sorumluluk siniri degisir; guvenlik kontrol matrisi bu modele gore ayrik tanimlanmalidir.

\section{Bulut Yerlisi (Cloud-Native) ve Konteyner Guvenligi}
\subsection{Docker, Kubernetes ve Kernel Izolasyon Teknolojileri}
Konteyner guvenligi icin imaj tarama, admission policy, runtime koruma ve namespace/NetworkPolicy kontrolleri uygulanmalidir.

\subsection{Kod Olarak Altyapi (IaC) Otomasyon Guvenligi}
IaC kodu (Terraform/CloudFormation vb.) merge oncesi guvenlik kurallarindan gecmeli, gizli anahtarlar koddan ayrik yonetilmelidir.
